% GDZ PPN236018647,native550lower–551;HKUST551–552.
\NumberedItem{73}
\indent Quod Iul. 27. inscrips[imus] errorem involvit: sed eo felicius nunc rem exhausimus, quoniam probare possumus nullum periodum esse posse numerum rationalem.
\par\noindent\hfill[1797] Aug. 1.

\medskip\begingroup\small
\indent Würden die Notizen Nr. 71 und 73 auf die Kreisteilung, nämlich auf die Einheitswurzeln vom Primzahlgrad bezogen, so wären sie unmittelbare Folgerungen aus der Aussage der Notiz Nr. 40 vom 9. Oktober 1796, d. h. aus der Irreduzibilität der Kreisteilungsgleichung; dann wäre aber bei Nr. 71 nichts Falsches. Wie S. \textsc{Gundelfinger} in einem Briefe an \textsc{Klein} vom 10. April 1903 bemerkt, wird aber \textsc{Gauss} in der Nr. 71 unter \(n\) jedenfalls eine beliebige zusammengesetzte Zahl verstanden haben, und es kämen demnach hier Sätze in Betracht, wie sie zuerst E. E. \textsc{Kummer}\footnote{E. E. \textsc{Kummer}, \textit{Theorie der idealen Primfaktoren der complexen Zahlen, welche aus den Wurzeln der Gleichung \(\omega^n=1\) gebildet sind, wenn \(n\) eine zusammengesetzte Zahl ist}, Abhandlungen der Akademie der Wissensch. zu Berlin (1856) 1857, Mathem. Abh., S. 1; siehe auch zwei Abhandlungen von L. \textsc{Fuchs} in \textsc{Crelles} Journal f. Mathematik \textbf{61}, 1863, S. 374 und \textbf{65}, 1866, S. 74, \textsc{Fuchs}' Werke I, S. 53 und S. 69.} aufgestellt hat. Bedeutet \(n\) eine zusammen-
\par\endgroup
\par\vfill\noindent X1.\hfill66

\clearpage\SourceLeaf{551}\ArithmeticHead{522}{L}{TAGEBUCH MIT ERLÄUTERUNGEN.}
\begingroup\small
\noindent gesetzte Zahl, so enthält die Nr. 71 in der Tat einen Irrtum, indem die Perioden \textls[120]{einer Gruppe} (nach \textsc{Kummer} a. a. O. § 2) dann einander gleich werden können, wenn sie verschwinden. Dies kann \textsc{Gauss}, wenn er, wie sehr wahrscheinlich, die Perioden ebenso wie \textsc{Kummer} definiert hat, nachträglich sehr wohl erkannt haben und die Nr. 73 bedeutet dann, wenn auch keinen Ersatz für die Aussage der Nr. 71, so doch einen erheblichen Fortschritt in der Erkenntnis dieser Verhältnisse. Unter einer \textls[120]{rationalen} Zahl wäre in der Nr. 73 eine \textls[120]{von Null verschiedene} zu denken.
\par\noindent\hfill\textsc{Bachmann.}
\par\endgroup

\NumberedItem{74}
\indent Quomodo periodorum numerum duplicando signa adornare oporteat.
\par\noindent\hfill[1797 Aug.]

\medskip\begingroup\small
\indent Hier ist wohl folgendes gemeint. Geht man von den \(e\) Perioden mit \(f\) Gliedern \(\eta_i\) (Voraussetzung ist ein gerades \(f\)) zu den \(2e\) Perioden mit \(\dfrac f2\) Gliedern \(\eta'_i\) über, so sind diese durch jene mittels quadratischer Gleichungen bestimmt und es fragt sich, welches Vorzeichen der bei der Auflösung dieser Gleichungen auftretenden Quadratwurzeln einer jeden der kleinern Perioden zukommt.
\par\noindent\hfill\textsc{Bachmann.}
\par\endgroup

\NumberedItem{75}
\indent Functionum primarum multitudinem per analysin simplicissimam erui.
\par\noindent\hfill[1797] Aug. 26.

\medskip\begingroup\small
\indent Der Anfang einer bezüglichen Tabelle ist bei S. 5 des \textsc{Leiste} aufgezeichnet; man vergl. die artt. 343, 344. der \textit{Analysis Residuorum}, Werke II, S. 220.
\par\noindent\hfill\textsc{Dedekind.}
\par\endgroup

\NumberedItem{76}
\indent \textls[120]{Theorema:} Si
\[
1+ax+bxx+\text{etc.}+mx^\mu
\]
\noindent est functio secundum modulum \(p\) prima, erit:
\[
d+x+x^p+x^{pp}+\text{etc.}+x^{p^{\mu-1}}
\]
\noindent per hanc f[un]ct[io]nem s[e]c[un]d[u]m hunc modulum divisibilis etc. etc.
\par\noindent\hfill[1797] Aug. 30.

\medskip\begingroup\small
\indent Vergl. den art. 356. der \textit{Analysis Residuorum}, Werke II, S. 227.
\par\noindent\hfill\textsc{Bachmann.}
\par\endgroup
